\documentclass{beamer}
\usetheme{Madrid}
\usecolortheme{whale}

\title{Top Quark Reconstruction with Machine Learning}
\subtitle{Revisiting Combinatorial Object Selection in ATLAS}
\author{GitHub Copilot}
\date{\today}

\begin{document}

\begin{frame}
  \titlepage
\end{frame}

\begin{frame}{Table of Contents}
  \tableofcontents
\end{frame}

\section{Concept}
\begin{frame}{The Problem: Top Quark Reconstruction}
  \begin{itemize}
    \item \textbf{Physics Process:} $t \to Wb \to q q' b$
    \item \textbf{The Challenge:} In realistic LHC events, we often have $N > 3$ jets due to radiation, pileup, and underlying events.
    \item \textbf{Combinatorics:} For $N$ jets, there are $\binom{N}{3}$ possible triplets.
    \item \textbf{Goal:} Identify the single triplet that corresponds to the true top decay among all candidates.
  \end{itemize}
  \vfill
  \centering
  \textit{Reconstruction = Structured Classification over Candidate Subsets}
\end{frame}

\begin{frame}{Feature Engineering}
  We use physics-motivated observables to distinguish signal from background:
  \begin{itemize}
    \item \textbf{Angular Geometry:} Three pairwise angular separations:
      \[ \Delta R_{ij} = \sqrt{(\eta_i - \eta_j)^2 + (\phi_i - \phi_j)^2} \]
    \item \textbf{Mass Structure:} Three pairwise invariant masses normalized to the total triplet mass:
      \[ \frac{m_{ij}}{m_{\text{triplet}}} \]
  \end{itemize}
  \vfill
  These six variables capture the internal kinematics of the decay while remaining robust and interpretable.
\end{frame}

\section{Implementation}
\begin{frame}{Package Architecture}
  The repository is structured as a modular ML pipeline (Source: \texttt{clean\_repo/docs/schema.md}):
  \begin{enumerate}
    \item \textbf{Stage 1: Dataset Build (\texttt{dataset\_build.py})}
      \begin{itemize}
        \item Processes ROOT ntuples.
        \item Labels triplets as truth or non-truth.
        \item Outputs \texttt{triplets\_raw.parquet}.
      \end{itemize}
    \item \textbf{Stage 2: Dataset Prepare (\texttt{dataset\_prepare.py})}
      \itemize{
        \item Splits data into Train, Val, and Test sets.
        \item Handles class balancing.
      }
    \item \textbf{Stage 3: Training (\texttt{train.py})}
      \itemize{
        \item Supports XGBoost (BDT) and TabPFN.
        \item Saves model artifacts and training reports.
      }
    \item \textbf{Stage 4: Inference (\texttt{infer.py})}
      \itemize{
        \item Evaluates model on test data.
        \item Ranks candidate triplets per event.
      }
  \end{enumerate}
\end{frame}

\begin{frame}{Key Dependencies}
  The implementation relies on modern data science libraries:
  \begin{itemize}
    \item \texttt{uproot} / \texttt{awkward}: For efficient ROOT file reading and vectorized jet manipulations.
    \item \texttt{pandas} / \texttt{pyarrow}: For triplet storage and Parquet management.
    \item \texttt{xgboost}: For the gradient boosted tree classifier.
    \item \texttt{matplotlib} / \texttt{plotly}: For automated diagnostic and performance plotting.
  \item \texttt{TabPFN}: For prior-data fitted network classification (optional).
  \end{itemize}
\end{frame}

\section{Operation}
\begin{frame}[fragile]{Standard Workflow (Makefile)}
  The recommended operation uses a \texttt{Makefile} for reproducibility:
  \begin{block}{Running the Pipeline}
\begin{verbatim}
# Set environment variables
export ROOT_INPUT=data/input.root

# Run the full XGBoost pipeline
make pipeline_xgb ROOT_INPUT=$ROOT_INPUT MAX_EVENTS=40000

# Run the TabPFN pipeline
make pipeline_tabpfn ROOT_INPUT=$ROOT_INPUT MAX_EVENTS=40000
\end{verbatim}
  \end{block}
  Outputs are organized under \texttt{artifacts/run\_\{N\}/}.
\end{frame}

\begin{frame}[fragile]{Direct CLI Interface}
  Individual stages can be run using the \texttt{triplet\_ml} module:
\begin{verbatim}
# Stage 1: Build
python -m triplet_ml dataset_build --inputs $ROOT_INPUT

# Stage 3: Train
python -m triplet_ml train \
    --model xgb \
    --xgb-config configs/xgb_hyperparameters.json \
    --train artifacts/train.parquet \
    --val artifacts/val.parquet
\end{verbatim}
\end{frame}

\begin{frame}{Monitoring and Visualization}
  The package automatically generates evaluation plots:
  \itemize{
    \item \textbf{Training logs:} ROC curves, loss over epochs, and feature importance.
    \item \textbf{Inference diagnostics:} Classification scores and invariant mass distributions of selected triplets.
    \item \textbf{Deployment:} Built-in support for deploying plots to public CFS directories via \texttt{--deploy}.
  }
\end{frame}

\begin{frame}{Summary}
  \begin{itemize}
    \item \textbf{Concept:} Frame top reconstruction as a supervised classification problem over triplets.
    \item \textbf{Implementation:} Modular, end-to-end Python pipeline using ROOT, Parquet, and XGBoost.
    \item \textbf{Operation:} Reproducible via \texttt{Makefile} or explicit CLI commands with structured artifact management.
  \end{itemize}
  \vfill
  \centering
  \Large \textbf{Questions?}
\end{frame}

\end{document}
